% Generated from a `canvaslms quizzes export` of the practice
% exam by mkexempel.py; do not edit by hand.
\begin{otherlanguage}{swedish}
\subsection{Question 1}
\label{ep3:q1}
\emph{Learning objective 1, 1 point, multi-answer, drawn from bank 1.1.}

Vilket eller vilka av följande namn är bra namn på en variabel som håller reda på antalet studenter? Kryssa i alla bra namn.

The options; every correct option is marked:
\begin{description}
  \item[{A}] \texttt{NUMBER\_\allowbreak{}OF\_\allowbreak{}STUDENTS}

    \emph{Feedback:} Det här är inget bra namn: \texttt{NUMBER\_\allowbreak{}OF\_\allowbreak{}STUDENTS} skrivs som en konstant. Versaler säger att värdet inte ska ändras, men det här värdet ändras.
  \item[{B \textit{(key)}}] \texttt{number\_\allowbreak{}of\_\allowbreak{}students}

    \emph{Feedback:} Det här är ett bra namn: \texttt{number\_\allowbreak{}of\_\allowbreak{}students} säger att det är ett antal studenter och är skrivet som en variabel.
  \item[{C}] \texttt{integer}

    \emph{Feedback:} Det här är inget bra namn: \texttt{integer} säger bara vilken typ värdet har, inte vad det betyder. Ett program har många värden av samma typ.
  \item[{D}] \texttt{many}

    \emph{Feedback:} Det här är inget bra namn: \texttt{many} säger att det är många, men inte av vad eller hur många.
\end{description}

The item's feedback:
\begin{description}
  \item[General feedback] \textbf{Tänk på:} ett bra namn säger vad värdet är, inte vilken typ det har, och en variabel skrivs med små bokstäver och understreck, eftersom värdet ändras. Frågan ger poäng bara om du kryssar i alla bra namn och inga andra, så pröva dem ett och ett.

\textbf{Prova:} terminalen hjälper inte, eftersom Python godtar alla namnen. Tänk dig i stället namnet på en rad långt ner i ett program, till exempel i en \texttt{print}: förstår du vad värdet är utan att se resten?

\textbf{I kursen:} vecka 38, föreläsningen \emph{Variabler och utskrifter}, avsnitten \emph{Variabelnamn och konstanter} (namnet är till för den som läser, inte för datorn), \emph{Konstanter} (versalerna är ett meddelande till läsaren) och \emph{PEP 8} (vad beskrivande namn och en gemensam stil ger den som läser). Bedömning: rättningsmatrisen i \emph{Laboration (1) funktioner, variabler och utskrifter}, kriterierna \emph{Beskrivande namn} och \emph{Konsekvent språk/typografi}.
  \item[Feedback on a right answer] Bara \texttt{number\_\allowbreak{}of\_\allowbreak{}students} är bra: det säger att det är ett antal studenter, och det är skrivet som en variabel ska skrivas enligt PEP 8.
  \item[Feedback on a wrong answer] Pröva varje namn mot två frågor: säger det vad värdet betyder, och är det skrivet som PEP 8 säger om just den här sortens namn? Ett namn som bara säger typen, eller bara en bokstav, klarar inte den första.
\end{description}

\subsection{Question 2}
\label{ep3:q2}
\emph{Learning objective 1, 1 point, choice, drawn from bank 1.2.}

Vilket av påståendena om koden nedan är sant?

\begin{minted}{python}
def count_passed(scores, pass_mark):
    """Returnerar hur många poäng i scores som är lägre än pass_mark."""
    count = 0
    for score in scores:
        if score >= pass_mark:
            count += 1
    return count
\end{minted}

The options; the key is marked:
\begin{description}
  \item[{A \textit{(key)}}] Docstringen säger lägre än \texttt{pass\_mark}, men koden räknar poäng som är minst \texttt{pass\_mark}.

    \emph{Feedback:} Rätt. Villkoret \texttt{score \textgreater{}=\allowbreak{} pass\_\allowbreak{}mark} räknar poäng som är minst \texttt{pass\_mark}, medan docstringen lovar dem som är lägre.
  \item[{B}] Docstringen borde inte nämna parametrarna, eftersom de redan syns på \texttt{def}-raden.

    \emph{Feedback:} Inte den här: kursens docstring nämner just parametrarna och säger vad som returneras, så att den som anropar slipper läsa koden.
  \item[{C}] Funktionen räknar de poäng som är lägre än \texttt{pass\_mark}, som docstringen säger.

    \emph{Feedback:} Inte den här: läs villkoret i \texttt{if}-satsen och prova \texttt{count\_\allowbreak{}passed(\allowbreak{}[40,\allowbreak{} 50,\allowbreak{} 60]\allowbreak{},\allowbreak{} 50)\allowbreak{}}. Räkna själv de poäng som är lägre än 50: får du samma svar?
  \item[{D}] Funktionen räknar de poäng som är högre än \texttt{pass\_mark}, och docstringen säger samma sak.

    \emph{Feedback:} Inte den här: läs villkoret och docstringen var för sig. \texttt{count\_\allowbreak{}passed(\allowbreak{}[40,\allowbreak{} 50,\allowbreak{} 60]\allowbreak{},\allowbreak{} 50)\allowbreak{}} räknar med \texttt{50}; skulle \texttt{50} räknas om villkoret vore högre än? Och vad lovar docstringen?
\end{description}

The item's feedback:
\begin{description}
  \item[General feedback] \textbf{Tänk på:} läs först docstringen och sedan koden, och jämför. Gör funktionen det docstringen säger?

\textbf{Prova:} kör \texttt{count\_\allowbreak{}passed(\allowbreak{}[40,\allowbreak{} 50,\allowbreak{} 60]\allowbreak{},\allowbreak{} 50)\allowbreak{}}. Räkna själv vad svaret borde bli enligt docstringen, och jämför.

\textbf{I kursen:} vecka 38, föreläsningen \emph{Funktioner}, avsnittet \emph{Docstrings} (docstringen är till för den som ska använda funktionen utan att läsa kroppen, och den visas av \texttt{help}). Bedömning: rättningsmatrisen i \emph{Laboration (1) funktioner, variabler och utskrifter}, kriteriet \emph{Kommentarer}.
  \item[Feedback on a right answer] Docstringen säger lägre än \texttt{pass\_mark}, men koden räknar poängen som är \texttt{\textgreater{}= pass\_mark}. \texttt{count\_\allowbreak{}passed(\allowbreak{}[40,\allowbreak{} 50,\allowbreak{} 60]\allowbreak{},\allowbreak{} 50)\allowbreak{}} ger \texttt{2}, inte \texttt{1}. En docstring som inte stämmer är värre än ingen, eftersom den som anropar litar på den.
  \item[Feedback on a wrong answer] Läs koden två gånger. Först som Python: går den att köra, och vad returnerar den? Sedan som den som ska anropa funktionen: får du veta vad den tar emot och vad den ger tillbaka utan att läsa kroppen, och säger kommentarerna något som koden inte redan säger?
\end{description}

\subsection{Question 3}
\label{ep3:q3}
\emph{Learning objective 1, 1 point, matching, drawn from bank 1.3.}

Para ihop varje namn med vad det är ett namn på, enligt PEP 8 och hur kursen namnger saker. Varje alternativ används en gång.

Match each name with a description; the key:
\begin{description}
  \item[{\texttt{enroll\_student}}] Ett funktionsnamn
  \item[{\texttt{NUMBER\_\allowbreak{}OF\_\allowbreak{}STUDENTS}}] Ett konstantnamn
  \item[{\texttt{students}}] Ett annat variabelnamn
  \item[{\texttt{j}}] Namnet på en räknare i en slinga
  \item[{\texttt{Student}}] Ett klassnamn
\end{description}
The item's feedback:
\begin{description}
  \item[General feedback] \textbf{Tänk på:} två saker avslöjar vad ett namn står för. Skrivsättet: versaler, stor begynnelsebokstav eller små bokstäver. Och ordet: ett verb gör något, ett substantiv är något. Frågan ger poäng bara om alla fem paren är rätt.

\textbf{Prova:} terminalen hjälper inte, eftersom Python godtar alla namnen till allt. Sortera i stället namnen efter skrivsätt först. Bland dem med små bokstäver: vilket är ett verb, och vilket är en enda bokstav?

\textbf{I kursen:} vecka 38, föreläsningen \emph{Variabler och utskrifter}, avsnitten \emph{Variabelnamn och konstanter} (namnet är till för den som läser, inte för datorn), \emph{Konstanter} (versalerna är ett meddelande till läsaren) och \emph{PEP 8} (vad beskrivande namn och en gemensam stil ger den som läser); \emph{Övning: Funktioner och variabler}, uppgiften \emph{Någon annans program} (varför funktioner gärna heter som ett verb och variabler som det de innehåller); vecka 41, föreläsningen \emph{Klasser och objekt}, avsnittet \emph{Klasser och objekt} (klassnamn). Bedömning: rättningsmatrisen i \emph{Laboration (1) funktioner, variabler och utskrifter}, kriterierna \emph{Beskrivande namn} och \emph{Konsekvent språk/typografi}.
  \item[Feedback on a right answer] \begin{itemize}
  \item \texttt{enroll\_student} är en funktion: namnet börjar med ett verb, eftersom en funktion gör något, och skrivs med små bokstäver och understreck.
  \item \texttt{NUMBER\_\allowbreak{}OF\_\allowbreak{}STUDENTS} är en konstant: versaler och understreck säger att värdet inte ska ändras.
  \item \texttt{j} är räknaren i en slinga: en enda bokstav säger inget om vad värdet betyder och passar därför bara för en räknare.
  \item \texttt{Student} är en klass: stor begynnelsebokstav i varje ord och inga understreck (CapWords) är PEP 8:s skrivsätt för klasser.
  \item \texttt{students} är den andra variabeln: ett substantiv säger vad värdet är, och det skrivs med små bokstäver och understreck.
\end{itemize}
  \item[Feedback on a wrong answer] Två förväxlingar är vanliga. Den första gäller de namn som skrivs med små bokstäver: ett namn som börjar med ett verb är en funktion, medan ett substantiv är en variabel. Den andra gäller räknaren och den andra variabeln: en enda bokstav säger ingenting om vad värdet betyder och passar därför bara för en räknare i en slinga, medan ett beskrivande namn hör till ett värde som betyder något. Versaler, och stor begynnelsebokstav i varje ord, är olika saker: det första är en konstant, det andra en klass.
\end{description}

\subsection{Question 4}
\label{ep3:q4}
\emph{Learning objective 2, 10 points, choice, drawn from bank 2.1.}

Funktionen \texttt{halvings} ska returnera hur många gånger ett heltal \texttt{number} som är 1 eller större kan heltalsdelas med 2 innan det har blivit 1. Till exempel ska \texttt{halvings(15)} ge \texttt{3} (15, 7, 3, 1).

\begin{minted}{python}
def halvings(number):
    """Returnerar antalet heltalsdelningar med 2 tills number är 1."""
    count = 0
    ____
        number //= 2
        count += 1
    return count
\end{minted}

Vilket villkor ska stå i luckan för att funktionen ska ge rätt svar för alla sådana tal?

The options; the key is marked:
\begin{description}
  \item[{A}] \leavevmode\par
\begin{minted}{python}
while number % 2 == 0 and number > 2:
\end{minted}


    \emph{Feedback:} Inte den här: slingan fortsätter bara så länge talet är jämnt och stannar vid första udda talet, och dessutom redan vid 2. Exemplet blir fel: \texttt{halvings(15)} ger \texttt{0}, eftersom 15 är udda.
  \item[{B \textit{(key)}}] \leavevmode\par
\begin{minted}{python}
while number > 1:
\end{minted}


    \emph{Feedback:} Rätt. Slingan delar så länge talet är större än 1, och heltalsdelningen gör att också udda tal kommer vidare.
  \item[{C}] \leavevmode\par
\begin{minted}{python}
while number > 2:
\end{minted}


    \emph{Feedback:} Inte den här: slingan stannar redan när talet är 2, så den sista delningen, från 2 till 1, räknas inte. Exemplet stämmer, eftersom 15 kommer till 3 och \texttt{3 // 2} är 1, men \texttt{halvings(16)} ger \texttt{3} (16, 8, 4, 2) och \texttt{halvings(2)} ger \texttt{0}.
  \item[{D}] \leavevmode\par
\begin{minted}{python}
while number // 2 > 1:
\end{minted}


    \emph{Feedback:} Inte den här: villkoret prövar halva talet, så slingan stannar redan när talet är 2 eller 3 och den sista delningen räknas inte. Exemplet blir fel: \texttt{halvings(15)} ger \texttt{2}.
\end{description}

The item's feedback:
\begin{description}
  \item[General feedback] \textbf{Tänk på:} två av villkoren ger 3 för 15. Fråga dig var slingan ska stanna: när talet är 2 finns en delning kvar, eftersom \texttt{2 // 2} är 1.

\textbf{Prova:} pröva fler tal än exemplet, också de minsta, och räkna delningarna för hand.

\textbf{I kursen:} vecka 40, föreläsningen \emph{Upprepningar}, avsnitten \emph{while: så länge villkoret gäller} och \emph{Siffersumman, tre gånger} (\texttt{//=}).
  \item[Feedback on a right answer] \texttt{number //= 2} är heltalsdelning, så 15 blir 7, 7 blir 3 och 3 blir 1: också udda tal kan delas vidare. Slingan ska fortsätta så länge talet är större än 1, också när det är 2, eftersom \texttt{2 // 2} är 1 och den delningen ska räknas.
  \item[Feedback on a wrong answer] Pröva villkoret du valde med fler tal än exemplet, också de minsta. Skriv upp talen efter varje delning med \texttt{//} och räkna dem.
\end{description}

\subsection{Question 5}
\label{ep3:q5}
\emph{Learning objective 2, 10 points, choice, drawn from bank 2.2.}

Funktionen ska returnera \texttt{\textquotedbl{}negativt\textquotedbl{}}, \texttt{\textquotedbl{}noll\textquotedbl{}} eller \texttt{\textquotedbl{}positivt\textquotedbl{}}.

\begin{minted}{python}
def describe(number):
    """Returnerar ett ord för tecknet på heltalet number."""
    if number < 0:
        word = "negativt"
    if number == 0:
        word = "noll"
    else:
        word = "positivt"
    return word
\end{minted}

\texttt{describe(-{}4)} ger \texttt{\textquotedbl{}positivt\textquotedbl{}}. Vilken förklaring stämmer?

The options; the key is marked:
\begin{description}
  \item[{A}] \texttt{else} hör till båda \texttt{if}-satserna och körs, eftersom ett av deras villkor är falskt.

    \emph{Feedback:} Inte den här: \texttt{else} hör bara till det \texttt{if} som står närmast ovanför. Det första \texttt{if} har ingen \texttt{else}.
  \item[{B \textit{(key)}}] Det andra \texttt{if} börjar ett nytt val. För −4 är \texttt{number == 0} falskt, så dess \texttt{else} ersätter \texttt{\textquotedbl{}negativt\textquotedbl{}}.

    \emph{Feedback:} Rätt. Två \texttt{if} efter varandra är två val, och båda görs; med \texttt{elif} hade det blivit ett val.
  \item[{C}] \texttt{word} kan bara få ett värde en gång, så den första tilldelningen i funktionen ignoreras.

    \emph{Feedback:} Inte den här: den första tilldelningen görs, men en senare tilldelning till samma variabel ersätter värdet.
  \item[{D}] Python prövar alla villkor och väljer den sista grenen vars villkor är sant, och \texttt{else} räknas som sant.

    \emph{Feedback:} Inte den här: \texttt{else} är inget villkor; den körs när villkoret i sitt eget \texttt{if} är falskt. Fråga dig hur många val koden gör.
\end{description}

The item's feedback:
\begin{description}
  \item[General feedback] \textbf{Tänk på:} hur många val gör koden, ett eller två? \texttt{else} hör till det \texttt{if} som står närmast ovanför.

\textbf{Prova:} anropa funktionen med ett negativt tal, med 0 och med ett positivt tal, och lägg till en utskrift av \texttt{word} efter varje tilldelning.

\textbf{I kursen:} vecka 39, föreläsningen \emph{Villkor och styrstrukturer}, avsnitten \emph{Förgrening: att välja väg} och \emph{Tre åldersgrupper} (två \texttt{if} mot \texttt{elif}).
  \item[Feedback on a right answer] Koden gör två val. I det första får \texttt{word} värdet \texttt{\textquotedbl{}negativt\textquotedbl{}}. Det andra \texttt{if} börjar ett nytt val: \texttt{-{}4 == 0} är falskt, så dess \texttt{else} ger \texttt{word} värdet \texttt{\textquotedbl{}positivt\textquotedbl{}}. Med \texttt{elif} i stället för det andra \texttt{if} hade det blivit ett enda val.
  \item[Feedback on a wrong answer] Följ \texttt{describe(-{}4)} rad för rad. Vilka tilldelningar till \texttt{word} görs, och i vilken ordning?
\end{description}

\subsection{Question 6}
\label{ep3:q6}
\emph{Learning objective 2, 10 points, ordering, drawn from bank 2.3.}

Ett program ska skriva ut:

\begin{minted}{text}
-
--
---
\end{minted}

Kodraderna visas utan indrag. Placera dem i den ordning som ger utskriften när varje rad får rätt indrag. Kom ihåg att en sträng gånger ett heltal upprepar strängen: \texttt{\textquotedbl{}-{}\textquotedbl{} * 3} är \texttt{\textquotedbl{}-{}-{}-{}\textquotedbl{}}.

The lines to order, as the item lists them:
\begin{description}
  \item[{A}] \leavevmode\par
\begin{minted}{python}
for j in range(i):
\end{minted}

  \item[{B}] \leavevmode\par
\begin{minted}{python}
for i in range(1, 3):
\end{minted}

  \item[{C}] \leavevmode\par
\begin{minted}{python}
print("-" * i, end="")
\end{minted}

  \item[{D}] \leavevmode\par
\begin{minted}{python}
print("-" * j)
\end{minted}

\end{description}
The key: the order B, A, C, D.

The item's feedback:
\begin{description}
  \item[General feedback] \textbf{Tänk på:} bestäm först vilka rader som måste komma före andra: en rad som använder en variabel måste stå efter raden som ger den ett värde. Fråga dig sedan för varje utskrift om den ska göras en gång per rad i utskriften (i den yttre slingan) eller en gång per varv i den inre. \texttt{end=\textquotedbl{}\textquotedbl{}} och \texttt{end=\textquotedbl{} \textquotedbl{}} gör att nästa utskrift fortsätter på samma rad; en \texttt{print} utan \texttt{end} avslutar raden.

\textbf{Prova:} skriv raderna i den ordning du tror, med indrag, och kör dem i terminalen. Jämför utskriften rad för rad med den i frågan. Stämmer den inte, så flytta en rad i taget och kör igen.

\textbf{I kursen:} vecka 40, föreläsningen \emph{Upprepningar}, avsnitten \emph{Vilka rader hör till slingan?} och \emph{Slingor i slingor} (multiplikationstabellen med \texttt{end=\textquotedbl{}\textquotedbl{}}); \emph{Övning: Upprepningar, listor och moduler}, uppgiften \emph{Multiplikationstabellen} (en sträng gånger ett tal).
  \item[Feedback on a right answer] \texttt{i} är 1 och 2, och den inre slingan går \texttt{i} varv, alltså tre varv sammanlagt: tre rader. Varje varv skriver \texttt{i} streck utan radbyte och sedan \texttt{j} streck, som avslutar raden: 1 + 0, 2 + 0 och 2 + 1.
  \item[Feedback on a wrong answer] Utskriften har tre rader. Hur många varv går den yttre slingan, och hur många går den inre i varje? Vilken av de två utskrifterna låter raden fortsätta?

Vanligt fel: att byta plats på de två utskrifterna. Då avslutas den första raden innan något skrivits på den: utskriften börjar med en tom rad.
\end{description}

\subsection{Question 7}
\label{ep3:q7}
\emph{Learning objective 3, 100 points, choice, drawn from bank 3.1.}

Funktionen \texttt{print\_row} ska skriva ut heltalen i en lista, alla från 0 till 999, på en rad. Varje tal ska stå längst till höger i en kolumn som är fyra tecken bred. För \texttt{[3, 7, 5]} ska utskriften bli:

\begin{minted}{python}
   3   7   5
def print_row(numbers):
    """Skriver ut talen i numbers på en rad, i kolumner fyra tecken breda."""
    for number in numbers:
        ...
    print()
\end{minted}

Vilket alternativ ska stå på raden med \texttt{...} för att funktionen ska fungera för alla sådana listor?

The options; the key is marked:
\begin{description}
  \item[{A \textit{(key)}}] \leavevmode\par
\begin{minted}{python}
print(f"{number:4d}", end="")
\end{minted}


    \emph{Feedback:} Rätt. Varje tal tar exakt fyra tecken, oavsett antalet siffror.
  \item[{B}] \leavevmode\par
\begin{minted}{python}
print(f"   {number}", end="")
\end{minted}


    \emph{Feedback:} Inte den här: tre mellanslag före ger fyra tecken bara för tal med en siffra. \texttt{[3, 14, 159, 2]} ger \texttt{   3   14   159   2}.
  \item[{C}] \leavevmode\par
\begin{minted}{python}
print(f"   {number}", end=" ")
\end{minted}


    \emph{Feedback:} Inte den här: tre mellanslag före ger fel bredd för tal med flera siffror, och \texttt{end=\textquotedbl{} \textquotedbl{}} lägger dessutom ett mellanslag efter varje fält, så redan \texttt{[3, 7, 5]} ger \texttt{   3    7    5 }.
\end{description}

The item's feedback:
\begin{description}
  \item[General feedback] \textbf{Tänk på:} två av alternativen ger exakt rätt utskrift för \texttt{[3, 7, 5]}. Frågan gäller alla tal från 0 till 999. Vad händer med kolumnerna när talen har olika många siffror?

\textbf{Prova:} skriv in funktionen med ett alternativ i taget och anropa den med tal av olika längd, till exempel \texttt{[3, 14, 159, 2]} och \texttt{[999, 0]}. Kör \texttt{print(\allowbreak{}f\textquotedbl{}\textbar{}\allowbreak{}\{3:4d\}\textbar{}\allowbreak{}\{159:4d\}\textbar{}\allowbreak{}\textquotedbl{})\allowbreak{}} för att se hur breda fälten är.

\textbf{I kursen:} vecka 40, föreläsningen \emph{Upprepningar}, avsnittet \emph{Slingor i slingor} (\texttt{:4d} och \texttt{end=\textquotedbl{}\textquotedbl{}}); vecka 41, föreläsningen \emph{Behållare: Uppslagslistor}, avsnittet \emph{Att skriva ut boken}.
  \item[Feedback on a right answer] \texttt{:4d} fyller ut varje tal till fyra tecken med mellanslag före, hur många siffror talet än har, och \texttt{end=\textquotedbl{}\textquotedbl{}} gör att inget kommer emellan fälten. \texttt{print()} sist avslutar raden.
  \item[Feedback on a wrong answer] Pröva med tal som har en, två och tre siffror. Står varje tal längst till höger i en kolumn som är exakt fyra tecken bred?
\end{description}

\subsection{Question 8}
\label{ep3:q8}
\emph{Learning objective 3, 100 points, choice, drawn from bank 3.2.}

Ett frågesportprogram vill att användaren skriver blåbärets latinska namn (\emph{Vaccinium myrtillus}).

Vilken rad ställer frågan tydligast för användaren?

The options; the key is marked:
\begin{description}
  \item[{A \textit{(key)}}] \leavevmode\par
\begin{minted}{python}
latin_name = input("Vad heter blåbär på latin? ")
\end{minted}


    \emph{Feedback:} Rätt. Frågan säger vilken växt och vilket namn.
  \item[{B}] \leavevmode\par
\begin{minted}{python}
latin_name = input("Vad heter blåbär? ")
\end{minted}


    \emph{Feedback:} Inte den här: frågan säger inte att det är det latinska namnet som efterfrågas, så de flesta svarar \texttt{blåbär} eller blir osäkra.
  \item[{C}] \leavevmode\par
\begin{minted}{python}
latin_name = input("Blåbär? ")
\end{minted}


    \emph{Feedback:} Inte den här: \texttt{Blåbär? } säger inte vad programmet vill veta om blåbär.
  \item[{D}] \leavevmode\par
\begin{minted}{python}
latin_name = input("Vilket latinskt namn har den här växten? ")
\end{minted}


    \emph{Feedback:} Inte den här: \texttt{den här växten} syftar på något användaren inte ser. Programmet vet vilken växt det gäller, men det gör inte användaren.
\end{description}

The item's feedback:
\begin{description}
  \item[General feedback] \textbf{Tänk på:} den som svarar ser bara frågan. Vet hen vilken växt det gäller och på vilket språk namnet ska skrivas?

\textbf{Prova:} läs frågorna högt för någon som inte har sett uppgiften. Vad svarar hen?

\textbf{I kursen:} vecka 39, föreläsningen \emph{Inmatning och felhantering}, avsnittet \emph{Att fråga så att någon förstår}.
  \item[Feedback on a right answer] Frågan säger både vilken växt (blåbär) och vilket namn som efterfrågas (det latinska), och lämnar ett mellanslag före svaret.
  \item[Feedback on a wrong answer] Pröva varje fråga: går det att svara rätt utan att ha sett programmet? Säger frågan både vilken växt och vilket språk?
\end{description}

\subsection{Question 9}
\label{ep3:q9}
\emph{Learning objective 3, 100 points, multi-answer, drawn from bank 3.3.}

Programmet nedan

\begin{minted}{python}
class FruitTree:
    """Ett fruktträd med sort och pris."""

    def __init__(self, fruit, price):
        """Skapar ett träd av sorten fruit som kostar price kronor."""
        self.fruit = fruit
        self.price = price


tree = FruitTree("äpple", 450)
print(tree)
\end{minted}

skriver ut något i stil med (talet är olika från körning till körning):

\begin{minted}{text}
<__main__.FruitTree object at 0x7f3b7ffefc10>
\end{minted}

Vilka påståenden förklarar utskriften? Kryssa i alla som stämmer.

The options; every correct option is marked:
\begin{description}
  \item[{A \textit{(key)}}] Klassen säger inte hur ett fruktträd ska bli text, så \texttt{print} gör som för alla klasser: klassens namn och en adress.

    \emph{Feedback:} Det här stämmer: utan en egen \texttt{\_\_str\_\_} får klassen den som alla klasser har, och den känner inte till attributen.
  \item[{B \textit{(key)}}] Talet efter \texttt{at} är var i minnet objektet ligger, och det kan bli ett annat nästa gång programmet körs.

    \emph{Feedback:} Det här stämmer: talet är objektets adress, och den bestäms när programmet körs.
  \item[{C \textit{(key)}}] Utskriften blir likadan om attributen görs privata, så länge klassen inte säger hur ett fruktträd ska bli text.

    \emph{Feedback:} Det här stämmer: \texttt{print} läser inte attributen, vare sig de är publika eller privata; det gör bara en \texttt{\_\_str\_\_} som returnerar dem.
  \item[{D}] Adressen visas för att objektet inte har gjorts om till en sträng med \texttt{str(tree)} före utskriften.

    \emph{Feedback:} Det här stämmer inte: \texttt{print} gör om objektet till text på samma sätt som \texttt{str(tree)}, så \texttt{print(str(tree))} ger samma utskrift.
\end{description}

The item's feedback:
\begin{description}
  \item[General feedback] \textbf{Tänk på:} frågan ger poäng bara om du har markerat alla påståenden som stämmer och inga andra, så pröva dem ett och ett. Vad gör \texttt{print} med ett objekt, och varifrån kommer texten när klassen inte själv har sagt hur den ska se ut?

\textbf{Prova:} pröva påståendena. Kör programmet två gånger och jämför talen. Lägg till en metod \texttt{\_\_str\_\_} som returnerar en f-sträng med attributen och se vad som ändras. Gör attributen privata och se om något ändras.

\textbf{I kursen:} vecka 41, föreläsningen \emph{Klasser och objekt}, avsnitten \emph{Med klasser} och \emph{Dundermetoder}; vecka 41, \emph{Övning: Klasser och objekt}, uppgiften \emph{Metoderna som redan fanns}.
  \item[Feedback on a right answer] Klassen har ingen \texttt{\_\_str\_\_}, så \texttt{print} använder den som alla klasser ärver. Den visar klassens namn och var objektet ligger i minnet, och den platsen kan vara en annan vid nästa körning. Om attributen är publika eller privata spelar ingen roll, och \texttt{str(tree)} ger samma text.
  \item[Feedback on a wrong answer] Skilj på vad som avgör utskriften (om klassen har sagt hur objektet ska bli text) och vad som inte gör det (om attributen är publika eller privata).
\end{description}

\subsection{Question 10}
\label{ep3:q10}
\emph{Learning objective 4, 1\,000 points, rich-fill-blank, drawn from bank 4.1.}

Programmet nedan ska lägga till dagens anteckning sist i dagboken \texttt{diary.txt}, utan att det som redan står där försvinner. Vad ska stå i rutan?

\begin{minted}[escapeinside=||]{python}
note = input("Dagens anteckning: ")
with open("diary.txt", |\fbox{\strut\enspace(1)\enspace}|) as file:
    file.write(f"{note}\n")
\end{minted}

The key for each blank:
\begin{description}
  \item[(1)] \texttt{\textquotedbl{}a\textquotedbl{}}, accepted as the regular expression \texttt{\textasciicircum{}\textbackslash{}s*\allowbreak{}[\textquotedbl{}\textquotesingle{}]\allowbreak{}(\allowbreak{}a\textbar{}\allowbreak{}at\textbar{}\allowbreak{}ta\textbar{}\allowbreak{}a\textbackslash{}+\allowbreak{}\textbar{}\allowbreak{}\textbackslash{}+\allowbreak{}a\textbar{}\allowbreak{}at\textbackslash{}+\allowbreak{}\textbar{}\allowbreak{}a\textbackslash{}+\allowbreak{}t\textbar{}\allowbreak{}ta\textbackslash{}+\allowbreak{}\textbar{}\allowbreak{}t\textbackslash{}+\allowbreak{}a\textbar{}\allowbreak{}\textbackslash{}+\allowbreak{}at\textbar{}\allowbreak{}\textbackslash{}+\allowbreak{}ta)\allowbreak{}[\textquotedbl{}\textquotesingle{}]\allowbreak{}\textbackslash{}s*\allowbreak{}\$}
\end{description}
The item's feedback:
\begin{description}
  \item[General feedback] \textbf{Tänk på:} det som redan står i dagboken ska stå kvar, och dagens anteckning ska hamna sist.

\textbf{Prova:} skapa \texttt{diary.txt} med en rad, kör programmet med ditt läge och skriv en anteckning. Läs sedan filen: båda raderna ska stå där, den nya sist.

\textbf{I kursen:} vecka 45, föreläsningen \emph{Arbeta med filer}, avsnittet \emph{Läsa och skriva} (sammanfattningen \emph{Fillägen}).
  \item[Feedback on a right answer] \texttt{\textquotedbl{}a\textquotedbl{}} (append) öppnar filen för att lägga till i slutet. Det som redan står där rörs inte, och \texttt{write} skriver efter det.
  \item[Feedback on a wrong answer] Vanliga fel: \texttt{\textquotedbl{}w\textquotedbl{}} tömmer dagboken innan något skrivs, så bara dagens anteckning blir kvar. \texttt{\textquotedbl{}r\textquotedbl{}} öppnar bara för läsning. \texttt{\textquotedbl{}x\textquotedbl{}} ger \texttt{FileExistsError}, eftersom dagboken redan finns.
\end{description}

\subsection{Question 11}
\label{ep3:q11}
\emph{Learning objective 4, 1\,000 points, ordering, drawn from bank 4.2.}

Placera raderna i rätt ordning, så att fåglarna i listan \texttt{birds} läggs till sist i loggfilen \texttt{birds.txt}, en fågel per rad. Det som redan står i filen ska stå kvar, och filen ska bara öppnas en gång. Raderna visas utan indrag; du bestämmer bara deras ordning. Anta att varje rad får det indrag som ordningen kräver.

The lines to order, as the item lists them:
\begin{description}
  \item[{A}] \leavevmode\par
\begin{minted}{text}
with open("birds.txt", "a") as file:
\end{minted}

  \item[{B}] \leavevmode\par
\begin{minted}{python}
for bird in birds:
\end{minted}

  \item[{C}] \leavevmode\par
\begin{minted}{text}
file.write(line)
\end{minted}

  \item[{D}] \leavevmode\par
\begin{minted}{python}
line = f"{bird}\n"
\end{minted}

\end{description}
The key: the order A, B, D, C.

The item's feedback:
\begin{description}
  \item[General feedback] \textbf{Tänk på:} det som redan står i loggfilen ska stå kvar, och filen ska öppnas en enda gång. Vilken rad ska då stå först?

\textbf{Prova:} skapa \texttt{birds.txt} med en rad, skriv in raderna i din ordning och kör. Läs filen: den gamla raden ska stå först och fåglarna efter.

\textbf{I kursen:} vecka 45, föreläsningen \emph{Arbeta med filer}, avsnitten \emph{Öppna och stänga} och \emph{Läsa och skriva} (sammanfattningen \emph{Fillägen}).
  \item[Feedback on a right answer] \texttt{with} står först och öppnar filen en gång, i läget \texttt{\textquotedbl{}a\textquotedbl{}}. Slingan skriver sedan en rad per fågel, efter det som redan stod i filen.
  \item[Feedback on a wrong answer] Läs kravet om att filen bara ska öppnas en gång.

Vanliga fel: \texttt{for} före \texttt{with} ger här faktiskt rätt innehåll, eftersom \texttt{\textquotedbl{}a\textquotedbl{}} inte tömmer filen, men då öppnas filen en gång per fågel. \texttt{file.write(line)} före raden som bygger \texttt{line} ger \texttt{NameError}.
\end{description}

\subsection{Question 12}
\label{ep3:q12}
\emph{Learning objective 4, 1\,000 points, choice, drawn from bank 4.3.}

Funktionen nedan sparar en lista i en fil.

\begin{minted}{python}
def save_items(filename, items):
    """Skriver strängarna i items till filen filename, en per rad."""
    with open(filename, "w") as file:
        for item in items:
            file.write(f"{item}\n")
\end{minted}

Filen \texttt{budget.txt} finns inte än. Efter de två anropen nedan står bara raden \texttt{el} i den. Vilket påstående förklarar varför?

\begin{minted}{text}
save_items("budget.txt", ["hyra", "mat"])
save_items("budget.txt", ["el"])
\end{minted}

The options; the key is marked:
\begin{description}
  \item[{A}] Filen stängs inte mellan anropen, så det andra anropet skriver i samma öppna fil, med början på första raden.

    \emph{Feedback:} Inte den här: \texttt{with} stänger filen när blocket lämnas, alltså redan när det första anropet är klart.
  \item[{B \textit{(key)}}] Båda anropen öppnar filen med \texttt{\textquotedbl{}w\textquotedbl{}}, som tömmer den vid öppnandet, så det andra börjar med en tom fil.

    \emph{Feedback:} Rätt. Det andra anropet börjar med en tom fil.
  \item[{C}] Varje \texttt{write} ersätter allt som redan står i filen, så bara den rad som skrevs sist blir kvar.

    \emph{Feedback:} Inte den här: efter det första anropet står båda raderna i filen, så två \texttt{write} i rad ersätter inte varandra.
  \item[{D}] \texttt{with} tar bort filen när blocket lämnas, så det andra anropet skapar en helt ny fil med samma namn.

    \emph{Feedback:} Inte den här: filen finns kvar efter det första anropet, med båda raderna. \texttt{with} stänger filen men tar inte bort den.
\end{description}

The item's feedback:
\begin{description}
  \item[General feedback] \textbf{Tänk på:} vad står i filen efter det första anropet, och vad händer med det när filen öppnas för andra gången?

\textbf{Prova:} att köra de två anropen visar bara det som redan står i frågan. Kör i stället bara det första anropet och läs filen. Pröva sedan att anropa funktionen på en fil som redan har en rad.

\textbf{I kursen:} vecka 45, föreläsningen \emph{Arbeta med filer}, avsnitten \emph{Öppna och stänga} och \emph{Läsa och skriva}; vecka 45, \emph{Övning: Filhantering}, avsnittet \emph{Frågan om filnamnet}.
  \item[Feedback on a right answer] Efter det första anropet står \texttt{hyra} och \texttt{mat} i filen. Det andra anropet öppnar filen med \texttt{\textquotedbl{}w\textquotedbl{}}, som tömmer den redan när den öppnas, och skriver sedan bara \texttt{el}.
  \item[Feedback on a wrong answer] Läs filen efter det första anropet. Står båda raderna där? Finns filen kvar? Vad måste då det andra anropet ha gjort?
\end{description}

\subsection{Question 13}
\label{ep3:q13}
\emph{Learning objective 5, 10\,000 points, multi-answer, drawn from bank 5.1.}

Ett program ska fråga efter användarens längd i hela centimeter och vikt i kilogram. Vikten kan ha decimaler, till exempel 72.5. Vilka av raderna frågar så att användaren förstår vad som ska matas in, \emph{och} ger ett värde av rätt typ? (Kryssa för alla som stämmer.)

The options; every correct option is marked:
\begin{description}
  \item[{A \textit{(key)}}] \leavevmode\par
\begin{minted}{python}
height_cm = int(input("Hur lång är du (cm)? "))
\end{minted}


    \emph{Feedback:} Det här stämmer: frågan säger vad och i vilken enhet, och hela centimeter passar ett \texttt{int}.
  \item[{B \textit{(key)}}] \leavevmode\par
\begin{minted}{python}
weight_kg = float(input("Hur mycket väger du (kg)? "))
\end{minted}


    \emph{Feedback:} Det här stämmer: frågan säger vad och i vilken enhet, och \texttt{float} klarar ett svar som \texttt{72.5}.
  \item[{C}] \leavevmode\par
\begin{minted}{python}
weight_kg = int(input("Hur mycket väger du (kg)? "))
\end{minted}


    \emph{Feedback:} Det här stämmer inte: frågan är tydlig, men \texttt{int(\textquotedbl{}72.5\textquotedbl{})} ger \texttt{ValueError}. En vikt kan ha decimaler, och då kraschar raden.
  \item[{D}] \leavevmode\par
\begin{minted}{python}
height_cm = int(input("x:"))
\end{minted}


    \emph{Feedback:} Det här stämmer inte: typen är rätt, men \texttt{x:} säger inte vad användaren ska svara eller i vilken enhet, och svaret hamnar tätt intill kolonet.
\end{description}

The item's feedback:
\begin{description}
  \item[General feedback] \textbf{Tänk på:} varje rad ska klara två krav: det ska framgå vad som ska skrivas och i vilken enhet, och värdet ska få en typ som passar svaret. Pröva varje rad mot båda. Frågan ger poäng bara om du markerar alla påståenden som stämmer och inga andra, så pröva dem ett och ett.

\textbf{Prova:} kör raderna i Python-terminalen. Svara \texttt{72.5} på viktfrågorna och \texttt{180} på längdfrågorna. Läs också frågan som om du inte hade sett koden: vet du vad du ska svara?

\textbf{I kursen:} vecka 39, föreläsningen \emph{Inmatning och felhantering}, avsnittet \emph{Att fråga så att någon förstår}.
  \item[Feedback on a right answer] En bra fråga säger vad som ska matas in och i vilken enhet, och slutar med ett mellanslag så att svaret inte klistras ihop med frågan. Längden i hela centimeter blir ett \texttt{int}, vikten med decimaler ett \texttt{float}. Kursens regel: fråga efter en sak i taget och skriv ut enheten i frågan.
  \item[Feedback on a wrong answer] En rad kan ge rätt typ och ändå fråga så att ingen förstår, eller fråga tydligt och ändå krascha på ett rimligt svar. Pröva varje rad mot båda kraven för sig.
\end{description}

\subsection{Question 14}
\label{ep3:q14}
\emph{Learning objective 5, 10\,000 points, choice, drawn from bank 5.2.}

Funktionen \texttt{input\_height} ska fråga tills användaren har svarat med ett tal och sedan returnera talet som ett \texttt{float}. Den får inte krascha, vad användaren än skriver. Vilken kropp ska stå i stället för \texttt{...}?

\begin{minted}{python}
def input_height():
    """Frågar tills svaret är ett tal och returnerar längden i meter."""
    ...
\end{minted}

The options; the key is marked:
\begin{description}
  \item[{A \textit{(key)}}] \leavevmode\par
\begin{minted}{python}
while True:
    try:
        height_m = float(input("Längd (m)? "))
        break
    except ValueError:
        print("Svara med ett tal.")
return height_m
\end{minted}


    \emph{Feedback:} Rätt. \texttt{break} står inne i \texttt{try}, efter omvandlingen, så slingan slutar bara när \texttt{float} har lyckats.
  \item[{B}] \leavevmode\par
\begin{minted}{python}
try:
    return float(input("Längd (m)? "))
except ValueError:
    print("Svara med ett tal.")
    input_height()
\end{minted}


    \emph{Feedback:} Inte den här: den frågar igen genom att anropa sig själv, men det svaret returneras aldrig. Efter ett felaktigt svar och sedan \texttt{1.75} blir returvärdet \texttt{None}.
  \item[{C}] \leavevmode\par
\begin{minted}{python}
while True:
    try:
        height_m = float(input("Längd (m)? "))
    except ValueError:
        print("Svara med ett tal.")
    break
return height_m
\end{minted}


    \emph{Feedback:} Inte den här: \texttt{break} står efter \texttt{try}, så slingan slutar efter första varvet vad som än hände. Efter ett felaktigt svar har \texttt{height\_m} inget värde, och \texttt{return} ger \texttt{UnboundLocalError}.
\end{description}

The item's feedback:
\begin{description}
  \item[General feedback] \textbf{Tänk på:} alla tre kropparna returnerar rätt om första svaret är ett tal. Frågan är vad som händer när det inte är det, och vad funktionen då returnerar.

\textbf{Prova:} skriv in funktionen med en kropp i taget och kör \texttt{print(\allowbreak{}input\_\allowbreak{}height(\allowbreak{})\allowbreak{})\allowbreak{}}. Svara inte med ett tal direkt: svara först med något som inte är ett tal, till exempel \texttt{1,75}, och sedan ett tal. Titta på vad som skrivs ut sist, inte bara på om det kraschar.

\textbf{I kursen:} vecka 39, föreläsningen \emph{Villkor och styrstrukturer}, avsnitten \emph{Att fråga om igen} och \emph{Det som kontrollen inte klarar}; vecka 40, föreläsningen \emph{Upprepningar}, avsnittet \emph{Inmatning som frågar om igen}.
  \item[Feedback on a right answer] Slingan står för omfrågningen och \texttt{try} för svar som inte går att omvandla. Var \texttt{break} står är det avgörande: inne i \texttt{try}, direkt efter omvandlingen, nås det bara när \texttt{float} har lyckats. Ett \texttt{ValueError} hoppar över \texttt{break}, skriver meddelandet och leder till ett nytt varv.
  \item[Feedback on a wrong answer] Felen syns bara efter ett felaktigt svar. Pröva varje kropp med ett felaktigt svar först och ett tal sedan, och se både om den kraschar och vad den returnerar.
\end{description}

\subsection{Question 15}
\label{ep3:q15}
\emph{Learning objective 5, 10\,000 points, choice, drawn from bank 5.3.}

Användaren gissar \texttt{7}, och programmet nedan skriver ut \texttt{Fel,\allowbreak{} försök igen en annan gång.\allowbreak{}} Vilken förklaring stämmer?

\begin{minted}{python}
SECRET = 7

guess = input("Gissa ett tal mellan 1 och 10: ")
if guess == SECRET:
    print("Rätt!")
else:
    print("Fel, försök igen en annan gång.")
\end{minted}

The options; the key is marked:
\begin{description}
  \item[{A \textit{(key)}}] \texttt{guess} är strängen \texttt{\textquotedbl{}7\textquotedbl{}}, och en sträng är aldrig lika med ett heltal, så villkoret är falskt för alla svar.

    \emph{Feedback:} Rätt. Värdena har olika typ, och då är de aldrig lika.
  \item[{B}] \texttt{input} tar med radbrytningen från Enter, så \texttt{guess} blir \texttt{\textquotedbl{}7\textbackslash{}n\textquotedbl{}}, och därför blir jämförelsen falsk.

    \emph{Feedback:} Inte den här: \texttt{input} tar inte med radbrytningen; svaret \texttt{7} blir \texttt{\textquotedbl{}7\textquotedbl{}}. Det är typen som gör jämförelsen falsk.
  \item[{C}] Python omvandlar svaret till ett tal i jämförelsen, men en konstant som \texttt{SECRET} kan bara jämföras med andra konstanter.

    \emph{Feedback:} Inte den här: Python omvandlar inget i en jämförelse, och en konstant är ett vanligt namn med ett värde. \texttt{int(\allowbreak{}\textquotedbl{}7\textquotedbl{})\allowbreak{} =\allowbreak{}=\allowbreak{} SECRET} är sant.
  \item[{D}] \texttt{==} frågar om två namn är samma variabel, inte om värdena är lika, och \texttt{guess} och \texttt{SECRET} är olika variabler.

    \emph{Feedback:} Inte den här: \texttt{==} jämför värden. \texttt{a = 7}, \texttt{b = 7} och \texttt{a == b} ger \texttt{True}, fast det är två olika namn.
\end{description}

The item's feedback:
\begin{description}
  \item[General feedback] \textbf{Tänk på:} körningen är given: en rätt gissning bedöms som fel. Fundera på vilka två värden \texttt{==} jämför, och vilka typer de har.

\textbf{Prova:} pröva förklaringarna: kör \texttt{input(\textquotedbl{}? \textquotedbl{})}, svara \texttt{7} och se exakt vad du får tillbaka; kör \texttt{\textquotedbl{}7\textquotedbl{} == 7} och \texttt{int(\textquotedbl{}7\textquotedbl{}) == 7}; skriv \texttt{a = 7}, \texttt{b = 7} och \texttt{a == b}.

\textbf{I kursen:} vecka 39, föreläsningen \emph{Inmatning och felhantering}, avsnitten \emph{Vad är det vi får tillbaka?} och \emph{Samma tecken, olika typ}; vecka 39, föreläsningen \emph{Villkor och styrstrukturer}, avsnittet \emph{Jämförelser ger sant eller falskt}.
  \item[Feedback on a right answer] \texttt{input} lämnar tillbaka strängen \texttt{\textquotedbl{}7\textquotedbl{}}, och \texttt{\textquotedbl{}7\textquotedbl{} == 7} är falskt: en sträng är aldrig lika med ett heltal. Python omvandlar inget i en jämförelse. Med \texttt{guess =\allowbreak{} int(\allowbreak{}input(\allowbreak{}.\allowbreak{}.\allowbreak{}.\allowbreak{})\allowbreak{})\allowbreak{}} jämförs två heltal, och gissningen \texttt{7} blir rätt.
  \item[Feedback on a wrong answer] Fråga dig vad \texttt{input} lämnar tillbaka, och om Python gör om något åt dig i en jämförelse. Jämför sedan resultatet av \texttt{==} med samma siffror i två olika typer.
\end{description}

\subsection{Question 16}
\label{ep3:q16}
\emph{Learning objective 6, 100\,000 points, choice, drawn from bank 6.1.}

Ett program har en klass \texttt{Meeting} där varje objekt representerar ett möte. Vad bör \texttt{is\_cancelled} vara i programmet? Det är \texttt{True} om just det mötet har ställts in och annars \texttt{False}; det sätts till \texttt{False} när mötet skapas och till \texttt{True} när mötet ställs in.

The options; the key is marked:
\begin{description}
  \item[{A \textit{(key)}}] Ett attribut

    \emph{Feedback:} Rätt. Det är ett värde som varje möte har och som sätts och ändras.
  \item[{B}] En metod

    \emph{Feedback:} Inte den här: namnet låter som en fråga, men beskrivningen säger att värdet \emph{sätts}, när mötet skapas och när det ställs in. Det som sätts och ändras är objektets tillstånd. Metoden är det som ställer in mötet och ändrar värdet.
  \item[{C}] En funktion utanför klassen

    \emph{Feedback:} Inte den här: en funktion gör något. Här beskrivs ett värde som hör till varje objekt.
  \item[{D}] En konstant

    \emph{Feedback:} Inte den här: en konstant är ett värde som är detsamma i hela programmet och aldrig ändras, som en momssats. Här skiljer sig värdet eller beteendet från objekt till objekt.
\end{description}

The item's feedback:
\begin{description}
  \item[General feedback] \textbf{Tänk på:} skilj på vad varje objekt \emph{har} och vad det \emph{gör}. Det ett objekt har, ett värde som kan vara olika för olika objekt, är ett attribut. Det objektet gör med sina egna värden är en metod. Det som inte hör till något särskilt objekt är en funktion utanför klassen, och ett värde som är detsamma i hela programmet och aldrig ändras är en konstant. Gå efter beskrivningen, inte efter hur namnet låter.

\textbf{Prova:} terminalen avgör inte frågan, eftersom alla fyra formerna går att skriva. Skissa i stället klassen: vad sätter \texttt{\_\_init\_\_} i \texttt{self}, och vilka metoder behöver \texttt{self} för att göra sitt jobb? Det som inte behöver \texttt{self} alls hör inte hemma i klassen.

\textbf{I kursen:} vecka 41, föreläsningen \emph{Klasser och objekt}, avsnitten \emph{Två exempel} (\emph{Utan klasser} och \emph{Med klasser}) och \emph{Metoder och attribut}; vecka 41, \emph{Övning: Klasser och objekt}, avsnittet \emph{Boken}.
  \item[Feedback on a right answer] \texttt{is\_cancelled} är ett värde, \texttt{True} eller \texttt{False}, som varje möte har. Det sätts i \texttt{\_\_init\_\_} och ändras när mötet ställs in, alltså är det objektets tillstånd: ett attribut. Att namnet låter som en fråga gör det inte till en metod.
  \item[Feedback on a wrong answer] Ställ två frågor om det som beskrivs. Är det ett värde eller något som görs? Hör det till ett visst objekt, eller är det detsamma för alla, eller hör det inte till objekten alls? Läs beskrivningen en gång till; ett namn kan låta som en sak och ändå vara en handling, eller tvärtom.
\end{description}

\subsection{Question 17}
\label{ep3:q17}
\emph{Learning objective 6, 100\,000 points, choice, drawn from bank 6.2.}

Betrakta följande kod:

\begin{minted}{python}
class Fruit:
    """En frukt med sort och pris."""

    def __init__(self, kind, price):
        """Skapar en frukt av sorten kind som kostar price kronor."""
        self.__kind = kind
        self.__price = price

    @property
    def kind(self):
        """Returnerar fruktens sort."""
        return self.__kind

    @property
    def price(self):
        """Returnerar fruktens pris i kronor."""
        return self.__price

    def __str__(self):
        """Skriver ut frukten med sort och pris."""
        print(f"{self.__kind}: {self.__price} kr")


fruit_salad = [Fruit("äpple", 5), Fruit("mango", 25), Fruit("apelsin", 10)]
for fruit in fruit_salad:
    print(fruit)
\end{minted}

När koden körs skrivs detta ut; den sista raden är felmeddelandet som programmet avbryts med:

\begin{minted}{text}
äpple: 5 kr
TypeError: __str__ returned non-string (type NoneType)
\end{minted}

Varför skrivs den första raden ut, och varför avbryts programmet sedan?

The options; the key is marked:
\begin{description}
  \item[{A \textit{(key)}}] \texttt{print(fruit)} anropar \texttt{\_\_str\_\_}, som skriver ut raden själv och sedan returnerar \texttt{None} i stället för en sträng.

    \emph{Feedback:} Rätt. Raden skrivs ut inne i \texttt{\_\_str\_\_}, och sedan får \texttt{print} \texttt{None} i stället för en sträng.
  \item[{B}] Bara den första frukten fick en \texttt{\_\_str\_\_}; de andra frukterna saknar metoden, så programmet avbryts vid den andra.

    \emph{Feedback:} Inte den här: metoderna står i klassen och gäller alla objekt av den. Alla tre frukterna har samma \texttt{\_\_str\_\_}, och programmet avbryts redan vid den första.
  \item[{C}] Egenskaperna \texttt{kind} och \texttt{price} returnerar \texttt{None} när de läses inifrån \texttt{\_\_str\_\_}, och \texttt{print} klarar inte \texttt{None}.

    \emph{Feedback:} Inte den här: egenskaperna returnerar värdena; \texttt{fruit\_\allowbreak{}salad[0]\allowbreak{}.\allowbreak{}kind} ger \texttt{\textquotedbl{}äpple\textquotedbl{}}. Raden som skrevs ut visar att \texttt{\_\_str\_\_} kunde läsa dem.
  \item[{D}] \texttt{print} skriver ut det \texttt{\_\_str\_\_} gav och skickar sedan sitt eget returvärde, \texttt{None}, tillbaka till \texttt{\_\_str\_\_}, som inte kan ta emot det.

    \emph{Feedback:} Inte den här: \texttt{print} anropar \texttt{\_\_str\_\_} en gång per objekt och skickar aldrig något tillbaka till metoden; en \texttt{\_\_str\_\_} som räknar sina anrop visar ett enda.
\end{description}

The item's feedback:
\begin{description}
  \item[General feedback] \textbf{Tänk på:} \texttt{print(fruit)} behöver göra om objektet till text och anropar då \texttt{\_\_str\_\_}. Vad gör den här \texttt{\_\_str\_\_}, och vad returnerar den?

\textbf{Prova:} utskriften står redan i frågan, så pröva förklaringarna. Anropa \texttt{fruit\_\allowbreak{}salad[0]\allowbreak{}.\allowbreak{}\_\allowbreak{}\_\allowbreak{}str\_\allowbreak{}\_\allowbreak{}(\allowbreak{})\allowbreak{}} direkt och skriv ut vad den returnerar. Läs också \texttt{fruit\_\allowbreak{}salad[1]\allowbreak{}.\allowbreak{}kind} för att se vad egenskaperna ger.

\textbf{I kursen:} vecka 41, föreläsningen \emph{Klasser och objekt}, avsnittet \emph{Med klasser}; vecka 41, \emph{Övning: Klasser och objekt}, avsnittet \emph{Boken}; föreläsningen 6/10, där en \texttt{\_\_str\_\_} som inte returnerade något gav \texttt{TypeError}.
  \item[Feedback on a right answer] \texttt{print(fruit)} anropar \texttt{fruit.\_\_str\_\_()} för att få en text att skriva ut. Den här \texttt{\_\_str\_\_} skriver själv ut raden, och det är därför \texttt{äpple: 5 kr} syns. Sedan tar metoden slut utan \texttt{return} och returnerar \texttt{None}, och \texttt{print} avbryter med \texttt{TypeError} eftersom den fick något som inte är en sträng. Med \texttt{return} i stället för \texttt{print} i \texttt{\_\_str\_\_} skrivs alla tre raderna ut.
  \item[Feedback on a wrong answer] Läs \texttt{\_\_str\_\_} rad för rad: vad skrivs ut, och vad returneras? Fråga dig sedan vad \texttt{print} gör med det metoden returnerar.
\end{description}

\subsection{Question 18}
\label{ep3:q18}
\emph{Learning objective 6, 100\,000 points, choice, drawn from bank 6.3.}

En spellista ska kunna räkna ut sin sammanlagda längd. Klasserna ser ut så här, men metoden saknas:

\begin{minted}{python}
class Song:
    """En låt med titel och längd i sekunder."""

    def __init__(self, title, length_s):
        """Skapar låten title som är length_s sekunder lång."""
        self.title = title
        self.length_s = length_s


class Playlist:
    """En spellista med namn och låtar."""

    def __init__(self, name):
        """Skapar en tom spellista med namnet name."""
        self.__name = name
        self.__songs = []

    def add(self, song):
        """Lägger till låten song sist i spellistan."""
        self.__songs.append(song)

    ...  # Här ska metoden total_length_s stå.
\end{minted}

Vilken version av \texttt{total\_length\_s} returnerar den sammanlagda längden för varje spellista?

The options; the key is marked:
\begin{description}
  \item[{A \textit{(key)}}] \leavevmode\par
\begin{minted}{python}
def total_length_s(self):
    """Returnerar spellistans sammanlagda längd i sekunder."""
    total = 0
    for song in self.__songs:
        total += song.length_s
    return total
\end{minted}


    \emph{Feedback:} Rätt. Den börjar på 0 och lägger till varje låt i listan en gång.
  \item[{B}] \leavevmode\par
\begin{minted}{python}
def total_length_s(self):
    """Returnerar spellistans sammanlagda längd i sekunder."""
    total = self.__songs[0].length_s
    for song in self.__songs[1:]:
        total += song.length_s
    return total
\end{minted}


    \emph{Feedback:} Inte den här: den ger rätt summa så länge spellistan har minst en låt, men en tom spellista ger \texttt{IndexError}, eftersom \texttt{self.\_\_songs[0]} inte finns.
  \item[{C}] \leavevmode\par
\begin{minted}{python}
def total_length_s(self):
    """Returnerar spellistans sammanlagda längd i sekunder."""
    total = 0
    counted = []
    for song in self.__songs:
        if song not in counted:
            total += song.length_s
            counted.append(song)
    return total
\end{minted}


    \emph{Feedback:} Inte den här: den hoppar över en låt som redan har räknats. Har samma låt lagts till två gånger spelas den två gånger, men räknas bara en gång.
  \item[{D}] \leavevmode\par
\begin{minted}{python}
def total_length_s(self):
    """Returnerar spellistans sammanlagda längd i sekunder."""
    total = self.__songs[0].length_s
    for song in self.__songs:
        total += song.length_s
    return total
\end{minted}


    \emph{Feedback:} Inte den här: den första låten räknas två gånger, en gång när \texttt{total} får sitt första värde och en gång i slingan. Två låtar på 231 och 168 sekunder ger 630 i stället för 399.
\end{description}

The item's feedback:
\begin{description}
  \item[General feedback] \textbf{Tänk på:} en spellista kan se ut på många sätt. Vilka spellistor kan programmet råka ut för, och klarar varje version dem alla?

\textbf{Prova:} skapa en spellista med två olika låtar och kontrollera summan; det räcker inte för att skilja versionerna åt. Pröva också en tom spellista, en med en enda låt och en där samma låt har lagts till två gånger, och räkna ut för hand vad svaret ska bli.

\textbf{I kursen:} vecka 44, anteckningarna \emph{Praktiska tillämpningar av klasser}, avsnitten \emph{Inköpslistor med klasser} och \emph{Banksystem med sammansatta klasser}; \emph{Övning: Fler behållare och mer om klasser}, uppgiften \emph{Banken}; Laboration 5; vecka 41, föreläsningen \emph{Klasser och objekt}, avsnittet \emph{Med klasser}; vecka 40, föreläsningen \emph{Upprepningar}, avsnittet \emph{Samma slinga över olika samlingar}.
  \item[Feedback on a right answer] Summan börjar på 0 och ökar med varje låts längd, en gång för varje låt i listan. Den ger 0 för en tom spellista och räknar samma låt två gånger om den har lagts till två gånger, precis som spellistan spelar den två gånger.
  \item[Feedback on a wrong answer] Pröva versionen du valde med fler spellistor än den självklara: en tom, en med en låt och en där samma låt finns två gånger. Vad ska summan bli, och vad ger versionen?
\end{description}

\subsection{Question 19}
\label{ep3:q19}
\emph{Learning objective 7, 1\,000\,000 points, multi-answer, drawn from bank 7.1.}

Vad är sant om att dela upp ett program i funktioner? Kryssa för alla påståenden som är sanna.

The options; every correct option is marked:
\begin{description}
  \item[{A}] Det blir svårare att hitta fel, eftersom programmet hoppar mellan funktionerna.

    \emph{Feedback:} Det här stämmer inte: felmeddelandet visar vilka anrop som ledde fram till felet och raden där det uppstod, så du hittar platsen även när programmet hoppar mellan funktioner. En funktion kan dessutom provas för sig, vilket gör felet lättare att ringa in.
  \item[{B \textit{(key)}}] Koden blir lättare att läsa när varje funktion gör en sak och har ett beskrivande namn.

    \emph{Feedback:} Det här stämmer: en funktion som gör en sak kan få ett namn som säger vad den gör, och då förstår läsaren ett anrop utan att läsa kroppen. Det är principen om ett ansvar (SRP) tillsammans med bra namn.
  \item[{C}] En funktion måste ha minst en parameter, annars har den inget att arbeta med.

    \emph{Feedback:} Det här stämmer inte: en funktion kan sakna parametrar, som \texttt{main()} i kursens program. Det funktionen arbetar med kan den skapa själv, till exempel genom att läsa in det med \texttt{input}.
  \item[{D}] Ett program måste innehålla egna funktioner för att gå att köra.

    \emph{Feedback:} Det här stämmer inte: Python kör filen uppifrån och ned, och en fil med bara \texttt{print(\textquotedbl{}Hej!\textquotedbl{})} är ett program som går att köra. Att kursens program alltid har funktioner och \texttt{main} är en regel för läsbarhet, inte något Python kräver.
  \item[{E \textit{(key)}}] Ska en beräkning ändras räcker det att ändra den i funktionen, eftersom koden bara står där.

    \emph{Feedback:} Det här stämmer: när koden bara står i funktionen räcker en ändring där, och alla anrop får den nya beräkningen. Med kopior måste varje kopia ändras, och glömmer du en räknar programmet på två sätt.
\end{description}

The item's feedback:
\begin{description}
  \item[General feedback] \textbf{Tänk på:} frågan är allt eller inget och ger poäng bara om du kryssar alla sanna påståenden och inga andra. Pröva dem ett i taget. Funktioner ändrar hur programmet är skrivet, för den som läser, ändrar, provar och återanvänder det, inte vad det gör när det körs. Har ett påstående ett skäl, ska både påståendet och skälet stämma.

\textbf{Prova:} terminalen kan pröva en del av påståendena: kör en fil som bara innehåller \texttt{print(\textquotedbl{}Hej!\textquotedbl{})}; skriv \texttt{def greet():} med raden \texttt{return \textquotedbl{}Hej\textquotedbl{}} och anropa \texttt{greet()}. De andra handlar om hur koden läses och ändras. Pröva dem på ett eget program: vad måste du ändra, och var, om programmet ska göra något nytt?

\textbf{I kursen:} vecka 38, föreläsningen \emph{Funktioner}, avsnitten \emph{Tre principer med namn}, \emph{Varför funktioner?}, \emph{Huvudprogrammet i en egen funktion} och \emph{Standardvärden}; vecka 39, föreläsningen \emph{Inmatning och felhantering}, avsnittet \emph{Att läsa ett felmeddelande}.
  \item[Feedback on a right answer] Det som stämmer är «Koden blir lättare att läsa när varje funktion gör en sak och har ett beskrivande namn» och «Ska en beräkning ändras räcker det att ändra den i funktionen, eftersom koden bara står där». Funktioner gör ett program lättare att läsa, ändra, prova och återanvända. Programmet gör samma sak som förut, och Python kräver dem inte.
  \item[Feedback on a wrong answer] Antingen har du kryssat något som inte stämmer, eller missat något som stämmer. Pröva varje påstående för sig: handlar det om hur koden läses, ändras, provas eller återanvänds, eller om vad som händer när programmet körs? Och håller skälet, om det har ett?
\end{description}

\subsection{Question 20}
\label{ep3:q20}
\emph{Learning objective 7, 1\,000\,000 points, multi-answer, drawn from bank 7.2.}

Du har skrivit funktionen nedan. Vilka av anropen är lämpliga, det vill säga tar hand om returvärdet (sparar det i en variabel eller använder det direkt) och ger inget felmeddelande? Varje anrop står för sig, direkt efter funktionen. Kryssa för alla som är lämpliga.

\begin{minted}{python}
def energy_cost(power_kw, price, hours=1):
    """Returnerar kostnaden för power_kw kW i hours timmar à price kr/kWh."""
    return power_kw * hours * price
\end{minted}

The options; every correct option is marked:
\begin{description}
  \item[{A}] \leavevmode\par
\begin{minted}{text}
energy_cost(2.5, 1.8, 3)
\end{minted}


    \emph{Feedback:} Det här anropet är inte lämpligt: antalet argument stämmer, men returvärdet varken sparas eller används, så det försvinner. I terminalens interaktiva läge, efter \texttt{\textgreater{}\textgreater{}\textgreater{}}, visas värdet ändå när anropet står ensamt på en rad, men i ett program syns ingenting.
  \item[{B \textit{(key)}}] \leavevmode\par
\begin{minted}{python}
cost = energy_cost(2.5, 1.8, 4)
\end{minted}


    \emph{Feedback:} Det här anropet är lämpligt: alla parametrar får ett argument, så standardvärdet används inte. Returvärdet sparas i \texttt{cost}.
  \item[{C \textit{(key)}}] \leavevmode\par
\begin{minted}{python}
cost = energy_cost(2.5, 1.8)
\end{minted}


    \emph{Feedback:} Det här anropet är lämpligt: \texttt{power\_kw} och \texttt{price} får var sitt argument, och \texttt{hours} får sitt standardvärde \texttt{1}. Returvärdet sparas i \texttt{cost}.
  \item[{D}] \leavevmode\par
\begin{minted}{python}
cost = energy_cost(2.5, 1.8, 3, 24)
\end{minted}


    \emph{Feedback:} Det här anropet är inte lämpligt: funktionen har tre parametrar men anropet ger fyra argument, så Python avbryter med \texttt{TypeError}.
  \item[{E}] \leavevmode\par
\begin{minted}{python}
cost = energy_cost(2.5)
\end{minted}


    \emph{Feedback:} Det här anropet är inte lämpligt: \texttt{price} har inget standardvärde och måste få ett argument, men anropet ger bara ett, så Python avbryter med \texttt{TypeError}.
\end{description}

The item's feedback:
\begin{description}
  \item[General feedback] \textbf{Tänk på:} frågan är allt eller inget, så kryssa alla lämpliga anrop och inga andra. Varje lämpligt anrop klarar två krav. Det första är rätt antal argument. \texttt{power\_kw} och \texttt{price} saknar standardvärde och måste alltid få ett argument; \texttt{hours} har ett standardvärde och får utelämnas. Ett anrop ska alltså ha minst 2 och högst 3 argument. Det andra är att värdet som funktionen returnerar sparas i en variabel eller används direkt, till exempel i en utskrift.

\textbf{Prova:} skriv in funktionen i terminalen och kör anropen ett i taget. Fel antal argument ger \texttt{TypeError}. Tänk på att terminalens interaktiva läge, efter \texttt{\textgreater{}\textgreater{}\textgreater{}}, visar värdet av ett anrop som står ensamt på en rad; i ett program försvinner det värdet utan att synas.

\textbf{I kursen:} vecka 38, föreläsningen \emph{Funktioner}, avsnitten \emph{Att lämna tillbaka ett värde}, \emph{Parametrar, argument och namn} och \emph{Standardvärden}.
  \item[Feedback on a right answer] De lämpliga anropen ger varje parameter utan standardvärde ett argument, inte fler argument än det finns parametrar, och tar hand om returvärdet: \texttt{cost =\allowbreak{} energy\_\allowbreak{}cost(\allowbreak{}2.\allowbreak{}5,\allowbreak{} 1.\allowbreak{}8,\allowbreak{} 4)\allowbreak{}} och \texttt{cost =\allowbreak{} energy\_\allowbreak{}cost(\allowbreak{}2.\allowbreak{}5,\allowbreak{} 1.\allowbreak{}8)\allowbreak{}}. Ett anrop är lämpligt först när båda kraven är uppfyllda.
  \item[Feedback on a wrong answer] Pröva varje anrop för sig mot de två kraven. Räkna argumenten och jämför med funktionshuvudet: en parameter utan standardvärde måste få ett argument, och fler argument än parametrar går inte. Titta sedan på vad som händer med värdet funktionen returnerar. Ett anrop kan klara det ena kravet men inte det andra.
\end{description}

\subsection{Question 21}
\label{ep3:q21}
\emph{Learning objective 7, 1\,000\,000 points, choice, drawn from bank 7.3.}

Vilket av programmen nedan är bäst uppdelat i funktioner? Alla tre programmen gör samma sak när de körs.

The options; the key is marked:
\begin{description}
  \item[{A \textit{(key)}}] \leavevmode\par
\begin{minted}{python}
"""Sammanfattar en veckas temperaturer."""


def mean(numbers):
    """Returnerar medelvärdet av talen i listan numbers."""
    total = 0
    for number in numbers:
        total += number
    return total / len(numbers)


def count_above(numbers, limit):
    """Returnerar hur många av talen i listan numbers som är över limit."""
    count = 0
    for number in numbers:
        if number > limit:
            count += 1
    return count


def main():
    """Skriver ut medeltemperaturen och antalet varma dagar."""
    temperatures = [12, -3, 7, 0, 15, 4, 9]
    print(f"Medeltemperaturen var {mean(temperatures):.1f} grader.")
    print(f"{count_above(temperatures, 5)} dagar var varmare än 5 grader.")


if __name__ == "__main__":
    main()
\end{minted}


    \emph{Feedback:} Rätt. Varje beräkning är en funktion med parametrar och returvärde, och \texttt{main} bara anropar och skriver ut.
  \item[{B}] \leavevmode\par
\begin{minted}{python}
"""Sammanfattar en veckas temperaturer."""


def mean():
    """Returnerar medelvärdet av temperaturerna."""
    total = 0
    for temperature in temperatures:
        total += temperature
    return total / len(temperatures)


def count_above():
    """Returnerar hur många av temperaturerna som är över 5 grader."""
    count = 0
    for temperature in temperatures:
        if temperature > 5:
            count += 1
    return count


temperatures = [12, -3, 7, 0, 15, 4, 9]
print(f"Medeltemperaturen var {mean():.1f} grader.")
print(f"{count_above()} dagar var varmare än 5 grader.")
\end{minted}


    \emph{Feedback:} Inte den här: funktionerna har inga parametrar. De läser listan \texttt{temperatures} från översta nivån, och gränsen 5 står inskriven i \texttt{count\_above}, så de fungerar bara för just den listan och den gränsen. Lägg huvudprogrammet i \texttt{main} och kör: det avbryts med \texttt{NameError: name \textquotesingle{}temperatures\textquotesingle{} is not defined}.
  \item[{C}] \leavevmode\par
\begin{minted}{python}
"""Sammanfattar en veckas temperaturer."""


def mean(numbers):
    """Returnerar medelvärdet av talen i listan numbers."""
    total = 0
    for number in numbers:
        total += number
    return total / len(numbers)


def main():
    """Skriver ut medeltemperaturen och antalet varma dagar."""
    temperatures = [12, -3, 7, 0, 15, 4, 9]
    print(f"Medeltemperaturen var {mean(temperatures):.1f} grader.")
    count = 0
    for temperature in temperatures:
        if temperature > 5:
            count += 1
    print(f"{count} dagar var varmare än 5 grader.")


if __name__ == "__main__":
    main()
\end{minted}


    \emph{Feedback:} Inte den här: medelvärdet är en funktion, men räkningen av varma dagar står direkt i \texttt{main}. Den går inte att använda igen, till exempel med en annan gräns, och \texttt{main} gör både sitt eget jobb och en beräkning.
\end{description}

The item's feedback:
\begin{description}
  \item[General feedback] \textbf{Tänk på:} alla tre programmen gör samma sak, så det är bara uppdelningen som skiljer dem. Gå igenom varje program med fyra frågor: Står huvudprogrammet i en egen funktion, \texttt{main}? Får varje funktion det den behöver som parametrar och lämnar sitt svar med \texttt{return}? Gör varje funktion en sak? Står någon kod två gånger? Fler funktioner är inte i sig bättre. Här: alla tre skriver ut samma två rader. Titta på var de två beräkningarna står och vad de får som parametrar: skulle du kunna använda dem för en annan lista eller en annan gräns?

\textbf{Prova:} kör programmen; alla tre skriver \texttt{Medeltemperaturen var 6.\allowbreak{}3 grader.\allowbreak{}} och \texttt{4 dagar var varmare än 5 grader.\allowbreak{}} Skriv sedan in funktionerna i terminalen och anropa dem med egna värden, till exempel \texttt{count\_\allowbreak{}above(\allowbreak{}[1,\allowbreak{} 8,\allowbreak{} 3]\allowbreak{},\allowbreak{} 2)\allowbreak{}}, som ska ge \texttt{2}. Vilket program låter dig göra det?

\textbf{I kursen:} vecka 38, föreläsningen \emph{Funktioner}, avsnitten \emph{Huvudprogrammet i en egen funktion}, \emph{Att dela upp ett program i funktioner} och \emph{Tre principer med namn}; vecka 40, föreläsningen \emph{Behållare: Listor}, avsnittet \emph{Att gå igenom en lista}.
  \item[Feedback on a right answer] Båda beräkningarna är funktioner som får listan (och gränsen) som parametrar och returnerar svaret. \texttt{main} skapar listan och skriver ut. Funktionerna fungerar för vilken lista och gräns som helst.
  \item[Feedback on a wrong answer] Ta programmet du valde och pröva det mot de fyra frågorna: huvudprogrammet i \texttt{main}, värden in som parametrar och ut med \texttt{return}, en uppgift per funktion, ingen kod två gånger. Ett program kan fungera och ändå bryta mot en av dem.
\end{description}
\end{otherlanguage}
